\chapter{Alcohols: Activating the OH Group}\label{ch:b1:alcohol-activation}

Alcohols are everywhere: made in tonnes from alkenes and from
fermentation, delivered by every Grignard addition of the previous
chapter. Yet an alcohol does not undergo substitution as a halogenoalkane
does: the OH group would have to leave as the hydroxide ion, a \omterm{def:b1:predominant-reaction:strong-acid}{strong
base} and a very poor \omterm{def:b1:electronic-effects:nucleophile}{leaving group} (\cref{ch:b1:nucleophilic-substitution}).
Chemists have three answers. Remove the proton of the OH instead, and the
\omterm{def:b1:alcohol-activation:alkoxide}{alkoxide} becomes a powerful \omterm{def:b1:electronic-effects:nucleophile}{nucleophile}. Add a proton, and the leaving
group becomes water. Or replace the hydrogen by a group that turns the
oxygen into part of an excellent \omterm{def:b1:electronic-effects:nucleophile}{leaving group}, without touching the
carbon. This chapter follows the three routes, and the choices of
mechanism and stereochemistry that come with them.

\begin{recall}
SN1, SN2 and their stereochemistry (\cref{ch:b1:nucleophilic-substitution});
E1, E2 and \omterm{prop:b1:elimination:zaitsev}{Zaitsev's rule} (\cref{ch:b1:elimination}); \omterm{def:b1:predominant-reaction:acidity-constant}{acidity constants}
and the predominant-reaction method (\cref{ch:b1:predominant-reaction});
the effect of alkyl groups on the acidity of alcohols
(\cref{exo:b1:electronic-effects:11}).
\end{recall}

\section{Alcohols as acids: alkoxides}

\begin{definition}[Alkoxide]\label{def:b1:alcohol-activation:alkoxide}
An \emph{alkoxide}\index{alkoxide} is the conjugate base \ce{RO-} of an
alcohol \ce{ROH}: methoxide \ce{CH3O-}, ethoxide \ce{CH3CH2O-},
tert-butoxide \ce{(CH3)3CO-}. Alkoxides are \omterm{def:b1:predominant-reaction:strong-acid}{strong bases} and good
\omterm{def:b1:electronic-effects:nucleophile}{nucleophiles}.
\end{definition}

\begin{proposition}[Acidity of alcohols]\label{prop:b1:alcohol-activation:alcohol-acidity}
Alcohols are weaker acids than water or about as weak: $\mathrm{p}K_a$
15.5 (methanol), 15.9 (ethanol), 17.1 (propan-2-ol), 19.2
(2-methylpropan-2-ol), against 14.0 for the couple \ce{H2O}/\ce{OH-}.
Hydroxide therefore deprotonates them only partly; a complete conversion
to the \omterm{def:b1:alcohol-activation:alkoxide}{alkoxide} needs a stronger base or an alkali metal.
\end{proposition}
% ledger: pka:methanol, pka:ethanol, pka:2-propanol, pka:tBuOH

\begin{proof}
The reaction \ce{ROH + OH- <=> RO- + H2O} has $K = 10^{14.0 -
\mathrm{p}K_a(\ce{ROH})}$, between $10^{-1.5}$ and $10^{-5}$: unfavourable.
Sodium metal reduces the proton, \ce{2ROH + 2Na -> 2RO- + 2Na+ + H2}, and
the hydride ion of sodium hydride takes it irreversibly,
\ce{ROH + H- -> RO- + H2}: the hydrogen escapes, and the reaction is
complete whatever its equilibrium constant.
\end{proof}

\begin{inthelab}[Sodium hydride]
Sodium hydride is sold as a grey powder dispersed in mineral oil, which
protects it from the moisture of the air. It reacts violently with water,
releasing flammable hydrogen; it is weighed quickly, added in portions
to the dry alcohol under an inert gas, and the bubbling of hydrogen shows
the progress of the deprotonation. Excess hydride is destroyed at the
end by slow addition of an alcohol, never of water.
\end{inthelab}

\section{The Williamson ether synthesis}

\begin{definition}[Williamson ether synthesis]\label{def:b1:alcohol-activation:williamson}
The \emph{Williamson ether synthesis}\index{Williamson ether synthesis}
prepares an ether $\mathrm{R{-}O{-}R'}$ by an SN2 reaction of an \omterm{def:b1:alcohol-activation:alkoxide}{alkoxide} \ce{RO-}
on a halogenoalkane (or a \omterm{def:b1:alcohol-activation:sulfonate}{sulfonate ester}) $\mathrm{R'{-}X}$.
\end{definition}

\begin{omfigure}
\setchemfig{atom sep=2.1em}
\schemestart
\chemfig{CH_3CH_2-@{o}O^{-}}\,\chemfig{Na^{+}}
\+
\chemfig{@{c}CH_3-[@{ci}]@{i}I}
\arrow{->}
\chemfig{CH_3CH_2-O-CH_3}
\+
\chemfig{Na^{+}}\,\chemfig{I^{-}}
\schemestop

\medskip
\chemmove{\draw[->, thick, shorten <=2pt, shorten >=2pt](o) .. controls +(80:8mm) and +(100:8mm) .. (c);
          \draw[->, thick, shorten <=2pt, shorten >=2pt](ci) .. controls +(90:5mm) and +(90:5mm) .. (i);}

{\small Williamson synthesis of methoxyethane: the ethoxide attacks the
carbon of iodomethane from the side opposite the iodine (SN2).}
\end{omfigure}

\begin{method}[Choosing the two partners]\label{met:b1:alcohol-activation:williamson-choice}
An ether $\mathrm{R{-}O{-}R'}$ can be cut on either side of the oxygen.
\begin{enumerate}
  \item Write the two possible pairs: \ce{RO-} with $\mathrm{R'X}$, and
        $\mathrm{R'O^-}$ with \ce{RX}.
  \item Keep the pair in which the halide is methyl or primary: SN2 needs
        an unhindered carbon.
  \item Reject a pair with a tertiary halide: the \omterm{def:b1:alcohol-activation:alkoxide}{alkoxide}, a \omterm{def:b1:predominant-reaction:strong-acid}{strong base},
        would eliminate (E2) instead.
  \item With a secondary halide, expect a mixture of ether and alkene.
\end{enumerate}
\end{method}

\begin{example}[An unsymmetrical ether]\label{ex:b1:alcohol-activation:williamson}
2-Methoxy-2-methylpropane (methyl tert-butyl ether): sodium tert-butoxide
with iodomethane gives it cleanly; sodium methoxide with the tertiary
bromide, 2-bromo-2-methyl\-propane, gives 2-methylpropene by E2. Bulky on the
\omterm{def:b1:alcohol-activation:alkoxide}{alkoxide} side, small on the halide side.
\end{example}

\section{Activating the OH group}

\begin{definition}[Sulfonate esters]\label{def:b1:alcohol-activation:sulfonate}
A \emph{sulfonate ester}\index{sulfonate ester} $\mathrm{R{-}O{-}SO_2{-}R'}$ is formed
from an alcohol and a sulfonyl chloride $\mathrm{R'SO_2Cl}$ in the presence of a
base (pyridine). The \emph{tosylate}\index{tosylate}, \ce{R-OTs}, comes
from 4-methylbenzenesulfonyl chloride (tosyl chloride); the
\emph{mesylate}\index{mesylate}, \ce{R-OMs}, from methanesulfonyl chloride.
The sulfonate anion $\mathrm{R'SO_3^-}$, the conjugate base of a \omterm{def:b1:predominant-reaction:strong-acid}{strong acid} and
stabilised by resonance over three oxygens, is an excellent \omterm{def:b1:electronic-effects:nucleophile}{leaving group}.
\end{definition}

\begin{omfigure}
\setchemfig{atom sep=2em}
\schemestart
\chemfig{R-O-H}
\+
\chemfig{Cl-S(=[2]O)(=[6]O)-[:0]*6(-=-(-CH_3)=-=)}
\arrow{->[pyridine]}
\chemfig{R-O-S(=[2]O)(=[6]O)-[:0]*6(-=-(-CH_3)=-=)}
\schemestop
\par\vspace{6pt}

{\small Tosylation of an alcohol. The O--H bond is replaced by O--S; the
C--O bond of the alcohol is not touched. Pyridine takes the HCl formed.}
\end{omfigure}

\begin{proposition}[Configuration through tosylation and substitution]\label{prop:b1:alcohol-activation:tosylation-retention}
Tosylation of an alcohol at a stereogenic carbon keeps its \omterm{def:b1:stereochemistry-in-depth:isomers}{configuration}
(retention); a subsequent SN2 substitution of the \omterm{def:b1:alcohol-activation:sulfonate}{tosylate} inverts it.
The two steps together replace OH by a \omterm{def:b1:electronic-effects:nucleophile}{nucleophile} with inversion.
\end{proposition}

\begin{proof}
In the tosylation, the bonds that change are O--H and S--Cl: the oxygen
of the alcohol attacks the sulfur and its hydrogen goes to the base. No
bond to the stereogenic carbon is broken, so the arrangement around it is
the same. In the SN2 step, the C--O bond breaks while the \omterm{def:b1:electronic-effects:nucleophile}{nucleophile}
bonds on the opposite side (\cref{prop:b1:nucleophilic-substitution:sn2-inversion}).
\end{proof}

\begin{omfigure}
\begin{tikzpicture}[font=\small, >={Stealth}]
  \node (a) at (0,0) {\chemfig{C(-[:150]H_3C)(<[:210]H)(<:[:250]C_2H_5)-OH}};
  \node at (0,-1.4) {\cip{S}-butan-2-ol};
  \draw[->, thick] (1.3,0) -- (2.5,0) node[midway, above] {TsCl} node[midway, below, font=\scriptsize] {retention};
  \node (b) at (3.9,0) {\chemfig{C(-[:150]H_3C)(<[:210]H)(<:[:250]C_2H_5)-OTs}};
  \node at (3.9,-1.4) {\cip{S} tosylate};
  \draw[->, thick] (5.3,0) -- (6.6,0) node[midway, above] {\ce{CN-}} node[midway, below, font=\scriptsize] {SN2, inversion};
  \node (c) at (8.3,0) {\chemfig{[:0]NC-C(-[:30]CH_3)(<[:-30]H)(<:[:-70]C_2H_5)}};
  \node at (8.3,-1.4) {\cip{R} nitrile};
\end{tikzpicture}

{\small From \cip{S}-butan-2-ol to 2-methylbutanenitrile: the tosylation
keeps the \omterm{def:b1:stereochemistry-in-depth:isomers}{configuration}, the SN2 step with cyanide inverts it. Overall:
OH replaced by CN with inversion.}
\end{omfigure}

The same goal can be reached in acid: the OH group, protonated, leaves as
water, \ce{R-OH2+ -> R+ + H2O} (SN1, E1) or under attack (SN2). Acid limits
the \omterm{def:b1:electronic-effects:nucleophile}{nucleophile}, though, to those that survive it --- halide ions, water,
alcohols --- and opens the door to \omterm{def:b1:electronic-effects:intermediates}{carbocations} and their side reactions.

\section{Conversion to halogenoalkanes}

\begin{proposition}[Alcohols and hydrogen halides]\label{prop:b1:alcohol-activation:hx-by-class}
Tertiary alcohols give the halogenoalkane with concentrated hydrochloric
acid at room temperature (SN1); primary alcohols need hydrobromic or
hydroiodic acid and heating (SN2 on the protonated alcohol); secondary
alcohols react at intermediate rates, often by both mechanisms.
\end{proposition}

\begin{proof}
Protonation turns OH into \ce{-OH2+}. For a tertiary alcohol, water leaves
easily, giving a stable \omterm{def:b1:electronic-effects:intermediates}{carbocation} captured by the halide (SN1): even the
modest \omterm{def:b1:electronic-effects:nucleophile}{nucleophile} \ce{Cl-} suffices. A primary \omterm{def:b1:electronic-effects:intermediates}{carbocation} is too
unstable; the halide must push water out from behind (SN2), which needs a
good \omterm{def:b1:electronic-effects:nucleophile}{nucleophile} (\ce{Br-}, \ce{I-}, better in protic media than \ce{Cl-},
\cref{prop:b1:electronic-effects:nucleophilicity-basicity}) and heat.
\end{proof}

\begin{example}[A classification test]\label{ex:b1:alcohol-activation:lucas}
A mixture of concentrated hydrochloric acid and zinc chloride (a \omterm{def:b1:lewis-resonance-vsepr:lewis-acid}{Lewis
acid} that helps the OH leave) turns a tertiary alcohol cloudy at once ---
the insoluble chloride separates --- a secondary alcohol within minutes,
and a primary alcohol hardly at all at room temperature: the order of
\omterm{def:b1:electronic-effects:intermediates}{carbocation} stability made visible in a test tube.
\end{example}

\section{Dehydration and ether formation in acid}

\begin{definition}[Dehydration of an alcohol]\label{def:b1:alcohol-activation:dehydration}
The \emph{dehydration}\index{dehydration} of an alcohol is the
elimination of water, catalysed by a \omterm{def:b1:predominant-reaction:strong-acid}{strong acid} (sulfuric or phosphoric
acid) and heat, giving an alkene: $\mathrm{R_2CH{-}CR'_2OH} \to \mathrm{R_2C{=}CR'_2} + \ce{H2O}$.
\end{definition}

\begin{omfigure}
\setchemfig{atom sep=2.1em}
\schemestart
\chemfig{C(-[2]CH_3)(-[:210]H_3C)(-[:-30]CH_2CH_3)-OH}
\arrow{->[\ce{H+}]}
\chemfig{C(-[2]CH_3)(-[:210]H_3C)(-[:-30]CH_2CH_3)-OH_2^{+}}
\arrow{->[\ce{-H2O}]}
\chemfig{C^{+}(-[2]CH_3)(-[:210]H_3C)(-[:-30]CH_2CH_3)}
\schemestop

\medskip
\schemestart
\chemfig{C^{+}(-[2]CH_3)(-[:210]H_3C)(-[:-30]CH_2CH_3)}
\arrow{->[\ce{-H+}][(major)]}
\chemfig{C(-[2]CH_3)(-[:210]H_3C)=[:-30]CHCH_3}
\schemestop

\medskip
{\small Acid-catalysed \omterm{def:b1:alcohol-activation:dehydration}{dehydration} of 2-methylbutan-2-ol (E1): protonation,
loss of water to the tertiary \omterm{def:b1:electronic-effects:intermediates}{carbocation}, loss of a proton. Removing the
proton from the \ce{CH2} gives the trisubstituted 2-methylbut-2-ene
(Zaitsev, major); from a methyl group, 2-methylbut-1-ene (minor).}
\end{omfigure}

\begin{proposition}[Ether or alkene]\label{prop:b1:alcohol-activation:ether-vs-alkene}
A primary alcohol heated with sulfuric acid gives mainly a symmetrical
ether at moderate temperature (one molecule substituting another by SN2)
and mainly an alkene at higher temperature; secondary and tertiary
alcohols give alkenes.
\end{proposition}

\begin{proof}
For a primary alcohol, the protonated alcohol can be attacked either at
carbon by a second alcohol molecule (SN2, giving \ce{R-O-R} after loss of
a proton) or at a $\beta$ hydrogen (E2). Elimination makes more molecules
than it consumes and is favoured as the temperature rises
(\cref{prop:b1:elimination:sn-vs-e}); substitution dominates when it is
lower. Secondary and tertiary alcohols ionise to \omterm{def:b1:electronic-effects:intermediates}{carbocations}, which lose a
proton more easily than they are captured by a weak \omterm{def:b1:electronic-effects:nucleophile}{nucleophile}, the more
so on heating; the water formed is also removed by distillation of the
volatile alkene.
\end{proof}

\section{Exercises}

\begin{exercise}[$\star$]\label{exo:b1:alcohol-activation:1}
Write the reactions of ethanol with sodium, with sodium hydride, and with
sodium hydroxide; compute the constant of the last ($\mathrm{p}K_a$ of
ethanol 15.9).
\end{exercise}

\begin{exercise}[$\star$]\label{exo:b1:alcohol-activation:2}
Give the products of: sodium methoxide with 1-bromopropane; sodium
ethoxide with iodomethane; sodium phenoxide with bromoethane.
\end{exercise}

\begin{exercise}[$\star$]\label{exo:b1:alcohol-activation:3}
Write the tosylation of propan-1-ol and the reaction of the \omterm{def:b1:alcohol-activation:sulfonate}{tosylate} with
sodium iodide. What is the advantage over the direct reaction of the
alcohol with sodium iodide?
\end{exercise}

\begin{exercise}[$\star$]\label{exo:b1:alcohol-activation:4}
Give the alkenes formed by acid \omterm{def:b1:alcohol-activation:dehydration}{dehydration} of butan-2-ol and of
2-methylpentan-2-ol, and the major one.
\end{exercise}

\begin{exercise}[$\star\star$]\label{exo:b1:alcohol-activation:5}
Propose a Williamson synthesis of 1-ethoxy-2-methylpropane and explain
your choice of partners. Why would the other choice give an alkene?
\end{exercise}

\begin{exercise}[$\star\star$]\label{exo:b1:alcohol-activation:6}
\cip{R}-octan-2-ol is converted into its \omterm{def:b1:alcohol-activation:sulfonate}{tosylate}, then treated with
sodium ethoxide in ethanol. Give the \omterm{def:b1:stereochemistry-in-depth:isomers}{configuration} of the ether. What
would be obtained by heating the same alcohol with sulfuric acid?
\end{exercise}

\begin{exercise}[$\star\star$]\label{exo:b1:alcohol-activation:7}
Write the mechanism of the reaction of 2-methylpropan-2-ol with
concentrated hydrochloric acid and explain why the reaction is fast at
room temperature.
\end{exercise}

\begin{exercise}[$\star\star$]\label{exo:b1:alcohol-activation:8}
Write the mechanism of the formation of ethoxyethane from ethanol in
sulfuric acid, and the competing elimination.
\end{exercise}

\begin{exercise}[$\star\star$]\label{exo:b1:alcohol-activation:9}
In the classification test of this chapter, which alcohol among
butan-1-ol, butan-2-ol and 2-methylpropan-2-ol turns cloudy first? Write
the product.
\end{exercise}

\begin{exercise}[$\star\star\star$]\label{exo:b1:alcohol-activation:10}
3,3-Dimethylbutan-2-ol, heated in acid, gives mainly
2,3-dimethylbut-2-ene, whose carbon skeleton differs from that of the
alcohol. Propose a mechanism with a 1,2-shift of a methyl group in the
\omterm{def:b1:electronic-effects:intermediates}{carbocation}, and justify the shift.
\end{exercise}

\begin{exercise}[$\star\star\star$]\label{exo:b1:alcohol-activation:11}
Design a synthesis of \cip{S}-2-methoxybutane from \cip{R}-butan-2-ol,
and another from \cip{S}-butan-2-ol, each keeping the \omterm{def:b1:stereochemistry-in-depth:isomers}{configuration} under
control.
\end{exercise}

\begin{exercise}[$\star\star\star$]\label{exo:b1:alcohol-activation:12}
Show that the equilibrium \ce{2CH3CH2OH <=> (CH3CH2)2O + H2O} can be
driven to the ether by removing water continuously, using the \omterm{def:b1:extent-q-and-k:quotient}{reaction
quotient} (\cref{ch:b1:extent-q-and-k}).
\end{exercise}

\section{Problem: Two Roads to an Ether}

\begin{problem}[{Weekend problem --- the two Williamson routes to methyl
tert-butyl ether, a stereochemically controlled ether from an optically
active alcohol, the dehydration of a tertiary alcohol, and the mass of
ether obtained}]
\label{pb:b1:alcohol-activation:1}
2-Methoxy-2-methylpropane (methyl tert-butyl ether, MTBE) was produced on a
large scale as a fuel additive. Molar masses (\unit{g/mol}): \ce{C4H10O}
74.0, \ce{C5H12O} 88.0, \ce{CH3I} 141.9.

\medskip
\noindent\textbf{Part I --- MTBE by the Williamson synthesis.}
\begin{enumerate}
  \item Write the two pairs of partners (\omterm{def:b1:alcohol-activation:alkoxide}{alkoxide} and halide) that could
        give MTBE.
  \item Which pair fails? Write the reaction that takes place instead.
  \item Write the successful reaction and its mechanism.
  \item How is sodium tert-butoxide prepared from 2-methylpropan-2-ol?
        Why not with sodium hydroxide?
  \item Why is the reaction run in the dry alcohol or an \omterm{def:b1:intermolecular-forces-solvents:solvent-classes}{aprotic solvent}?
  \item Is the product \omterm{def:b1:stereochemistry-in-depth:stereocentre}{chiral}?
\end{enumerate}

\medskip
\noindent\textbf{Part II --- An optically active ether.}
\begin{enumerate}[resume]
  \item \cip{R}-Butan-2-ol is converted into its \omterm{def:b1:alcohol-activation:sulfonate}{tosylate}. Write the
        reaction and give the \omterm{def:b1:stereochemistry-in-depth:isomers}{configuration} of the \omterm{def:b1:alcohol-activation:sulfonate}{tosylate}.
  \item The \omterm{def:b1:alcohol-activation:sulfonate}{tosylate} reacts with sodium methoxide. Write the product and
        its \omterm{def:b1:stereochemistry-in-depth:isomers}{configuration}.
  \item Which \omterm{def:b1:stereochemistry-in-depth:isomers}{configuration} of the alcohol gives \cip{R}-2-methoxybutane
        by the same route?
  \item Another route: the \omterm{def:b1:alcohol-activation:alkoxide}{alkoxide} of \cip{R}-butan-2-ol with
        iodomethane. \omterm{def:b1:stereochemistry-in-depth:isomers}{Configuration} of the product? Why?
  \item Which side reaction competes in question 8, and why is it
        limited here?
  \item Why would heating \cip{R}-butan-2-ol with methanol in sulfuric
        acid give a nearly racemic ether, if any?
\end{enumerate}

\medskip
\noindent\textbf{Part III --- \omterm{def:b1:alcohol-activation:dehydration}{Dehydration} of a tertiary alcohol.}
\begin{enumerate}[resume]
  \item Write the mechanism of the acid \omterm{def:b1:alcohol-activation:dehydration}{dehydration} of
        2-methylbutan-2-ol.
  \item Give the two alkenes and predict the major one.
  \item Why is no ether formed from this alcohol?
  \item Why is 2-methylbut-2-ene distilled out as it forms?
  \item Draw the \omterm{def:b1:elementary-steps:energy-profile}{energy profile} of the E1 step and mark the
        \omterm{def:b1:elementary-steps:rds}{rate-determining step}.
  \item Could a primary alcohol be dehydrated by the same mechanism?
\end{enumerate}

\medskip
\noindent\textbf{Part IV --- Masses.}
\begin{enumerate}[resume]
  \item \qty{10.0}{g} of 2-methylpropan-2-ol are converted into the
        \omterm{def:b1:alcohol-activation:alkoxide}{alkoxide}. Compute the amount.
  \item What mass of iodomethane is needed for a \qty{10}{\percent} excess?
  \item Compute the theoretical mass of MTBE.
  \item Compute the mass obtained for a \qty{75}{\percent} \omterm{def:b1:extent-q-and-k:fractional-extent}{yield}.
\end{enumerate}
\end{problem}
